\section{Completely satisfiable systems}\label{sec: p}

When studying systems where it is possible
to simultaneously satisfy all equations
the situation changes and let us start with
the positive results.

\subsection{Finding a good assignment}

Suppose we have an
equation of the form
$P(x)=1$ and that
this equation implies the affine
condition $A(x)=1$.  Then,
as the system is satisfiable,
we can use this equation to eliminate
one variable from the system, preserving the
degrees of all equations.  This is done
by taking some variable $x_i$ that appears
in $A$ and replacing it by $x_i+A(x)+1$ (note that
this function does not depend on $x_i$ as the two
occurrences of this variable cancel).
This substitution preserves the satisfiability of the
system and the process stops only 
when none of the current equations implies an affine condition.

Using \expref{Theorem}{basic theorem} we see
that when this process ends each equation
is satisfied by at least a fraction
$2^{1-d}-2^{1-2d}$ of the inputs and thus the following
theorem should not come as a surprise.

\begin{theorem}\label{thm upper sat}
There is a polynomial time algorithm that,
given a system of $m$ simultaneously satisfiable equations
of degree $d$ over $\GF[2]$, finds an
assignment that satisfies at least
$\lceil (2^{1-d}-2^{1-2d})m\rceil$ equations.
\end{theorem}

\begin{proof}
There are two points in the outlined argument
that require closer inspection.  The
first is the question of how to actually determine whether a 
polynomial equation implies an affine condition
and the second is to make sure that once
the process of finding implied affine conditions
has ended we can indeed deterministically find a
solution that satisfies the expected number of
equations.
Let us first address the issue of determining
whether a given equation implies an affine
condition.

Suppose $P(x)=1$ implies $A(x)=1$
for some unknown affine function $A$.
Let us assume that $x_1$ appears in $A$ with
a nonzero coefficient.  
We may write
$$
P(x)=P_0(x)+P_1(x)x_1$$
where neither $P_0$ nor $P_1$ depends on $x_1$.  We want
to prove that $P(x)=A(x)P_1(x)$ and towards this consider
\begin{eqnarray}\label{eq pq}
Q(x)=P(x)+A(x)P_1(x)\,.
\end{eqnarray}
As $x_1$ appears with coefficient one in $A$ it follows
that $Q$ does not depend on $x_1$ and let us assume that
$Q$ is not identically 0.  Choose any values
for $x_2, x_3 \ldots x_n$, to make $Q(x)=1$
and set $x_1$ to make $A(x)=0$.  It follows
from \eqref{eq pq} that $P(x)=1$ 
and thus we have found a counterexample 
to the assumed implication.  We can hence conclude
that $Q\equiv 0$ and we have 
$$
P(x)=A(x) P_1(x)\,.
$$

We claim furthermore that this procedure is entirely
efficient.  Namely given $P$ and the identity of
one variable occurring in $A$, $P_1$ is uniquely
defined.  Once $P_1$ is determined the rest
of the coefficients of $A$ can be found by
solving a linear system of equations.
As there are only $n$ candidates for a variable
in $A$ and solving a linear system of equations
can be done in polynomial time we conclude that the entire
process of finding implied affine conditions
can be done in polynomial time.

Once this process halts we need to find an assignment that satisfies the
expected number of equations.  This could
possibly be done in a manner similar
to how it was done in the proof of Theorem~3.1
but for the general case we do not know how to, in deterministic 
polynomial time,
find a suitable bound for the probability that an equation
is satisfied.  It is true that the bound of \expref{Theorem}{basic theorem}
can be calculated deterministically, but as we are not able to
prove it by induction, setting only one variable at the time,
it is not clear how to use it to guide an algorithm.
We turn to a different approach and we need the following
result by Viola~\cite{viola09}.

\begin{theorem}[\cite{viola09}]\label{thm:viola}
For any constant $d$ and $\epsilon > 0$ there exists
a polynomial time pseudo-random generator $G$ with
seed length $O(d \log n + d 2^d \log (1/\epsilon))$
and output in $\{ 0,1\}^n$ such that
$$
|\Pr[P(x)=0]- \Pr[P(G(y))=0]| \leq \epsilon
$$
for any polynomial $P$ of degree $d$.  Here $x$ is 
a uniformly random string in $\{ 0,1\}^n$ and
$y$ is a uniformly random seed.
\end{theorem}

This theorem implies that if we are willing
to sacrifice a small $\epsilon$ it is enough
to consider the candidate assignments 
$G(y)$ for all possible seeds $y$ and taking the assignment
from this set that satisfies the %
maximum
number of equations.  Let us be more precise.

We know that the expected number of satisfied
equations when we consider a uniformly random
input is $(2^{1-d}- 2^{1-2d}) m$ and, by
\expref{Theorem}{thm:viola}, the same number, when
we consider a random output from $G$, is
at least $(2^{1-d}- 2^{1-2d}-\epsilon) m$.
Setting $\epsilon= 2^{-2d} m^{-1}$ and observing that
the maximum number of equations satisfied is
an integer we see that this number is at least
$$
\lceil (2^{1-d}- 2^{1-2d}) m-2^{-2d} \rceil \,.
$$
Finally, as $(2^{1-d}- 2^{1-2d}) m$ is an integer
multiple of $2^{1-2d}$, we have that
$$
\lceil (2^{1-d}- 2^{1-2d}) m-2^{-2d} \rceil 
=\lceil (2^{1-d}- 2^{1-2d}) m \rceil\,,
$$
and the proof is complete.
\end{proof}

Note that the generator of Viola is only needed to
derandomize the algorithm and had we been content
with a randomized algorithms we could simply replace
the last step by random sampling, obtaining a much
more efficient algorithm.
